\documentclass[../../main.tex]{subfiles}% 注意这里的文件路径不能用 ./main.tex ,否则用latexmk编译子文件会报错
\graphicspath{{\subfix{./image/}}{\subfix{../../image/}}} % 指定图片引用路径,后续可以直接使用图片文件名
% 如果图片存放在上级目录,则使用 ../../image/ 作为备用路径

% 例如:
% \begin{figure}[H]
% \centering
% \includegraphics[width=0.3\textwidth]{图.png}
% \caption{}
% \label{figure:图}
% \end{figure}
% 注意:上述\label{}一定要放在\caption{}之后,否则引用图片序号会只会显示??.

\begin{document}

\section{上半平面的Dirichlet问题}

\begin{theorem}
\label{thm:dirichlet}
设 \(u\) 是定义在实轴上的逐段连续函数，如果 \(\lim\limits_{t\to+\infty}u(t)\) 和 \(\lim\limits_{t\to-\infty}u(t)\) 都存在且有限，那么以 \(u(t)\) 为边值的上半平面的 Dirichlet 问题的解可以写成
\begin{align}\label{eq:dirichlet_half_plane_1}
u(z_0)=\frac{1}{\pi} \int_{-\infty}^{\infty} \frac{y_0}{(t-x_0)^2+y_0^2}u(t)\mathrm{d}t, \qquad \operatorname{Im} z_0>0. 
\end{align}
也可以写成
\begin{align}\label{eq:dirichlet_half_plane_2}
u(z_0)=\operatorname{Re} \frac{1}{\pi i} \int_{-\infty}^{\infty} \frac{u(t)}{t-z_0}\mathrm{d}t.
\end{align}
\end{theorem}

\begin{proof}
在上半平面中任取一点 \(z_0=x_0+iy_0\)，作分式线性变换
\begin{align*}
w=\varphi(z)=\frac{z-z_0}{z-\overline{z}_0},
\end{align*}
它把上半平面一一地映为 \(|w|<1\)，把实轴一一地映为单位圆周 \(|w|=1\)。通过这一映射，\(u(t)\) 变成了 \(|\zeta|=1\) 上的函数 \(u_1(\zeta)=u(\varphi^{-1}(\zeta))\)，它在 \(|\zeta|=1\) 上除了有限个第一类间断点外处处连续。
根据\hyperref[theorem:poisson_theta]{单位圆盘上 Poisson 公式}，可以得到
\begin{align}
u_1(0)=\frac{1}{2\pi} \int_{0}^{2\pi} u_1(e^{i\varphi})\mathrm{d}\varphi. \label{eq:poisson_disk}
\end{align}
因为 \(u_1(0)=u(z_0)\)，\(u_1(e^{i\varphi})=u(t)\)，并且
\begin{align*}
e^{i\varphi}=\frac{t-z_0}{t-\overline{z}_0}.
\end{align*}
上式两边取对数后求微分，即得
\begin{align}
\mathrm{d}\varphi=\frac{2y_0\mathrm{d}t}{(t-x_0)^2+y_0^2}, \qquad z_0=x_0+iy_0. \label{eq:diff_phi}
\end{align}
把 \eqref{eq:poisson_disk} 式和 \eqref{eq:diff_phi} 式代入，即得 \(u(z)\) 的表达式
\begin{align*}
u(z_0)=\frac{1}{\pi} \int_{-\infty}^{\infty} \frac{y_0}{(t-x_0)^2+y_0^2}u(t)\mathrm{d}t, \qquad \operatorname{Im} z_0>0. 
\end{align*}
由于
\begin{align*}
\operatorname{Re}\frac{1}{i(t-z_0)}=\frac{y_0}{(t-x_0)^2+y_0^2},
\end{align*}
故上式也可改写为
\begin{align*}
u(z_0)=\operatorname{Re} \frac{1}{\pi i} \int_{-\infty}^{\infty} \frac{u(t)}{t-z_0}\mathrm{d}t.
\end{align*}\end{proof}

\begin{example}
设 \(w=f(z)\) 是把上半平面 \(D\) 一一地映为多角形区域 \(G\)，且在 \(\overline{D}\) 上连续的双全纯函数。如果 \(G\) 在其顶点 \(w_k\) 处的顶角为 \(\alpha_k\pi\)（\(0<\alpha_k<2\)），实轴上与 \(w_k\) 对应的点为 \(a_k\)（\(k=1,\cdots,n\)），\(-\infty<a_1<\cdots<a_n<\infty\)，则要证明
\begin{align}
f(z)=C\int_{z_0}^{z}(z-a_1)^{\alpha_1-1}\cdots(z-a_n)^{\alpha_n-1}\mathrm{d}z+C_1. \label{eq:schwarz_christoffel}
\end{align}
\end{example}

\begin{proof}
记 \(v(z)=\arg f'(z)\)，它是上半平面中的调和函数。考虑它在实轴上满足的条件。当 \(z\in(a_k,a_{k+1})\)（\(k=1,\cdots,n-1\)）时，\(f(z)\) 在线段 \(w_kw_{k+1}\) 上。根据导数辐角的几何意义，此时有
\begin{align*}
v(z)=\arg f'(z)=\arg(w_{k+1}-w_k), \qquad k=1,\cdots,n-1.
\end{align*}
若记 \(a_0=-\infty\)，\(a_{n+1}=\infty\)，则当 \(z\in(a_0,a_1)\) 或 \(z\in(a_n,a_{n+1})\) 时，\(f(z)\) 在线段 \(w_nw_1\) 上，所以此时有
\begin{align*}
v(z)=\arg f'(z)=\arg(w_1-w_n).
\end{align*}
如果记
\begin{align}
\theta_k=
\begin{cases}
\arg(w_{k+1}-w_k), & k=1,\cdots,n-1; \\
\arg(w_1-w_n), & k=0,n.
\end{cases} \label{eq:theta_k_def}
\end{align}
那么 \(v\) 在实轴上应满足条件
\begin{align*}
v(t)=\theta_k, \qquad a_k<t<a_{k+1}, \quad k=0,1,\cdots,n.
\end{align*}
根据上半平面 Dirichlet 问题解的公式 \eqref{eq:dirichlet_half_plane_2}，有
\begin{align*}
v(z)=\operatorname{Re} \frac{1}{\pi i}\int_{-\infty}^{\infty}\frac{v(t)}{t-z}\mathrm{d}t =\sum_{k=0}^{n} \operatorname{Re} \frac{\theta_k}{\pi i}\ln\frac{z-a_{k+1}}{z-a_k} =\frac{1}{\pi}\sum_{k=0}^{n} \theta_k \arg\frac{z-a_{k+1}}{z-a_k}.
\end{align*}
容易知道，当 \(z\) 在上半平面时，有
\begin{align}
\arg(z-a_0)=0, \qquad \arg(z-a_{n+1})=\pi. \label{eq:arg_boundary}
\end{align}
记
\begin{align*}
S_k=\sum_{j=0}^{k} \arg\frac{z-a_{j+1}}{z-a_j}=\arg(z-a_{k+1}).
\end{align*}
应用 Abel 变换和 \eqref{eq:arg_boundary} 式，有
\begin{align*}
v(z)=\frac{1}{\pi}\sum_{k=1}^{n}(\theta_{k-1}-\theta_k)\arg(z-a_k)+\theta_n.
\end{align*}
但从 \eqref{eq:theta_k_def} 式容易直接验证
\begin{align*}
\theta_{k-1}-\theta_k=(\alpha_k-1)\pi, \qquad k=1,\cdots,n.
\end{align*}
因而有
\begin{align*}
\arg f'(z)=v(z)=\sum_{k=1}^{n}(\alpha_k-1)\arg(z-a_k)+\theta_n.
\end{align*}
现在令
\begin{align*}
F(z)=\sum_{k=1}^{n}(\alpha_k-1)\ln(z-a_k)+i\theta_n,
\end{align*}
它和 \(\ln f'(z)\) 有相同的虚部，所以
\begin{align*}
\ln f'(z)=\ln C+i\theta_n+\sum_{k=1}^{n}(\alpha_k-1)\ln(z-a_k).
\end{align*}
由此即得
\begin{align*}
f'(z)=C\prod_{k=1}^{n}(z-a_k)^{\alpha_k-1}.
\end{align*}
进而求出积分，即得到公式 \eqref{eq:schwarz_christoffel}。\end{proof}




\end{document}